% ---
% title: A fair or a two-headed coin
% tags: bayes, odds
% ---

\begin{problem}
	A gambler has two coins in his pocket: a fair coin and a two-headed coin
	(heads on both sides). He picks one of them at random and flips it.
	(a) The coin shows heads. What is the probability that it is the fair
	coin? (b) He flips the same coin again, and it shows heads again. What is
	the probability now that it is the fair coin? (c) He flips the same coin a
	third time, and it shows tails. What is the probability now that it is the
	fair coin? Solve (a) and (b) twice: with Bayes' formula and with odds.
\end{problem}

\begin{solution}
	\idea Let $F$ be the event that he picked the fair coin; then $F^c$ is the
	event that he picked the two-headed coin. Before any flip,
	$P(F)=P(F^c)=\frac12$. Each flip is new information, and Bayes' formula
	turns the old probability of $F$ into a new one.

	\step{(a) Bayes' formula} Let $H_1$ be the event ``the first flip shows
	heads''. We know the probability of heads for each coin: the fair coin
	shows heads with probability $P(H_1\mid F)=\frac12$, the two-headed coin
	always shows heads, so $P(H_1\mid F^c)=1$. The events $F$ and $F^c$ split
	the sample space into two parts, so by Bayes' formula
	\[
		P(F\mid H_1)
		=\frac{P(H_1\mid F)\,P(F)}{P(H_1\mid F)\,P(F)+P(H_1\mid F^c)\,P(F^c)}
		=\frac{\frac12\cdot\frac12}{\frac12\cdot\frac12+1\cdot\frac12}
		=\frac{\frac14}{\frac34}
		=\frac13.
	\]
	Heads is more typical of the two-headed coin, so seeing heads makes the
	fair coin less likely: its probability drops from $\frac12$ to $\frac13$.

	\step{(a) Odds} The odds form of Bayes' formula says
	\[
		\underbrace{\frac{P(F\mid H_1)}{P(F^c\mid H_1)}}_{\text{new odds}}
		=\underbrace{\frac{P(F)}{P(F^c)}}_{\text{prior odds}}\cdot
		\underbrace{\frac{P(H_1\mid F)}{P(H_1\mid F^c)}}_{\text{likelihood ratio}}.
	\]
	The prior odds of the fair coin are $\frac{1/2}{1/2}=1$ (one to one). The
	likelihood ratio of heads is $\frac{1/2}{1}=\frac12$. So the new odds are
	$1\cdot\frac12=\frac12$, that is, $1:2$ for the fair coin. Odds of $1:2$
	mean that the fair coin has $1$ chance out of $1+2=3$:
	$P(F\mid H_1)=\frac{1/2}{1+1/2}=\frac13$, the same as before.

	\step{(b) Bayes' formula} Let $H_2$ be the event ``the second flip shows
	heads''. Given the coin, the two flips are independent. The fair coin shows
	two heads with probability $P(H_1H_2\mid F)=\frac12\cdot\frac12=\frac14$;
	the two-headed coin does so with probability $P(H_1H_2\mid F^c)=1$. Bayes'
	formula with the original prior $\frac12$ gives
	\[
		P(F\mid H_1H_2)
		=\frac{\frac14\cdot\frac12}{\frac14\cdot\frac12+1\cdot\frac12}
		=\frac{\frac18}{\frac58}
		=\frac15.
	\]
	The same result comes out step by step: after the first heads the
	probability of the fair coin is $\frac13$ (from (a)). Use this as the new
	prior for the second flip:
	\[
		\frac{\frac12\cdot\frac13}{\frac12\cdot\frac13+1\cdot\frac23}
		=\frac{\frac16}{\frac56}
		=\frac15.
	\]

	\step{(b) Odds} With odds the update is just one more multiplication. The
	odds after the first heads were $\frac12$; the second heads has the same
	likelihood ratio $\frac12$, so the new odds are
	$\frac12\cdot\frac12=\frac14$, that is, $1:4$. Hence
	$P(F\mid H_1H_2)=\frac{1/4}{1+1/4}=\frac15$.

	\step{(c)} Let $T_3$ be the event ``the third flip shows tails''. The
	two-headed coin can never show tails: $P(T_3\mid F^c\,H_1H_2)=0$. So after
	a tail the two-headed coin is impossible, and the coin must be the fair
	one. With Bayes' formula (prior $\frac15$ from (b), fair coin gives tails
	with probability $\frac12$):
	\[
		P(F\mid H_1H_2T_3)
		=\frac{\frac12\cdot\frac15}{\frac12\cdot\frac15+0\cdot\frac45}
		=\frac{\frac1{10}}{\frac1{10}}
		=1.
	\]

	\begin{answer}
		(a) $\dfrac13$ \quad (b) $\dfrac15$ \quad (c) $1$
	\end{answer}

	\begin{remark}
		Each observation updates the probability: every heads multiplies the
		odds of the fair coin by $\frac12$, so after $n$ heads in a row the odds
		are $\frac1{2^n}$, and the fair coin becomes less and less likely, but
		never impossible. One decisive observation settles the question: a
		single tail rules out the two-headed coin completely. With odds, the
		likelihood ratio of tails is $\frac{1/2}{0}$, which is not a number;
		this is the odds way of saying that the evidence is conclusive.
	\end{remark}
\end{solution}
